% GENERATED FILE. Edit canonical lessons, then run npm run generate:books.
% canonical-source-hash: fnv1a64:8d1edfb577fe54f5
% canonical-lessons: IT-C531-bucato, IT-C531-appendere, IT-C531-stendino, IT-C531-molletta, IT-C531-ferro-da-stiro

\chapter{Il bucato}
\label{ch:it-il-bucato}
\hlchaptermodality{531}

\begin{chapteropening}
\textbf{By the end of this chapter:} \emph{I can talk about doing the laundry.}

Complete the last lesson of chapter 531: I can talk about doing the laundry.
\end{chapteropening}

\section[il bucato]{il bucato — laundry, washing}

\label{lesson:IT-C531-bucato}



\begin{quote}
Before the new one: say the Italian for to post, then the Italian for follower.
\end{quote}

\subsection*{You'll want to know: il bucato}
\textbf{il bucato} — \textquotedblleft{}laundry, washing\textquotedblright{}.

\begin{quote}
Il sabato faccio il bucato.
\end{quote}

\begin{quote}
\emph{On Saturdays I do the washing.}
\end{quote}

\emph{Fare il bucato} — to do the washing.

\begin{cousinweb}
The origin of \emph{bucato} is not settled; it is usually traced to a Germanic verb for soaking clothes in lye.
\end{cousinweb}

\begin{grammarlens}[title={What you've built}]
One more word for telling, explaining and describing at B1.
\end{grammarlens}

\subsection*{Your turn}
\begin{itemize}
\raggedright
  \item \emph{Say it:} \emph{il bucato}
  \item \emph{Say it:} \emph{il bucato}, once more
  \item \emph{Say it:} a sentence of your own with \emph{il bucato}
  \item \emph{Recall:} say the Italian for to post, then the Italian for follower, then say \emph{il bucato} again
\end{itemize}

\begin{tcolorbox}[breakable,colback=teal!4,colframe=teal!35!black,title={Before you move on}]
What does il bucato mean? (\textquotedblleft{}Laundry, washing\textquotedblright{}.) Say it once more.
\end{tcolorbox}
\section[appendere]{appendere — to hang (up)}

\label{lesson:IT-C531-appendere}



\begin{quote}
Before the new one: say the Italian for follower, then the Italian for laundry.
\end{quote}

\subsection*{You'll want to know: appendere}
\textbf{appendere} — \textquotedblleft{}to hang (up)\textquotedblright{}.

\begin{quote}
Appendo il cappotto all'ingresso.
\end{quote}

\begin{quote}
\emph{I hang my coat up in the hall.}
\end{quote}

The past participle is \emph{appeso}: \emph{il quadro è appeso al muro}.

\begin{cousinweb}
\emph{appendere} is inherited from Latin \emph{appendere} \textquotedblleft{}to hang up, to weigh out\textquotedblright{} (EtymDB), from \emph{pendere} \textquotedblleft{}to hang\textquotedblright{}.
\end{cousinweb}

\begin{grammarlens}[title={What you've built}]
One more word for telling, explaining and describing at B1.
\end{grammarlens}

\subsection*{Your turn}
\begin{itemize}
\raggedright
  \item \emph{Say it:} \emph{appendere}
  \item \emph{Say it:} \emph{appendere}, once more
  \item \emph{Say it:} a sentence of your own with \emph{appendere}
  \item \emph{Recall:} say the Italian for follower, then the Italian for laundry, then say \emph{appendere} again
\end{itemize}

\begin{tcolorbox}[breakable,colback=teal!4,colframe=teal!35!black,title={Before you move on}]
What does appendere mean? (\textquotedblleft{}To hang (up)\textquotedblright{}.) Say it once more.
\end{tcolorbox}
\section[lo stendino]{lo stendino — clothes horse, drying rack}

\label{lesson:IT-C531-stendino}



\begin{quote}
Before the new one: say the Italian for laundry, then the Italian for to hang.
\end{quote}

\subsection*{You'll want to know: lo stendino}
\textbf{lo stendino} — \textquotedblleft{}clothes horse, drying rack\textquotedblright{}.

\begin{quote}
Metto lo stendino sul balcone.
\end{quote}

\begin{quote}
\emph{I put the clothes horse out on the balcony.}
\end{quote}

Built on \emph{stendere} \textquotedblleft{}to spread out, to hang out\textquotedblright{}.

\begin{cousinweb}
\emph{stendino} is built on \emph{stendere} \textquotedblleft{}to spread out\textquotedblright{}, inherited from Latin \emph{extendere} \textquotedblleft{}to stretch out\textquotedblright{} (Lewis \& Short), from \emph{tendere}.
\end{cousinweb}

\begin{grammarlens}[title={What you've built}]
One more word for telling, explaining and describing at B1.
\end{grammarlens}

\subsection*{Your turn}
\begin{itemize}
\raggedright
  \item \emph{Say it:} \emph{lo stendino}
  \item \emph{Say it:} \emph{lo stendino}, once more
  \item \emph{Say it:} a sentence of your own with \emph{lo stendino}
  \item \emph{Recall:} say the Italian for laundry, then the Italian for to hang, then say \emph{lo stendino} again
\end{itemize}

\begin{tcolorbox}[breakable,colback=teal!4,colframe=teal!35!black,title={Before you move on}]
What does lo stendino mean? (\textquotedblleft{}Clothes horse, drying rack\textquotedblright{}.) Say it once more.
\end{tcolorbox}
\section[la molletta]{la molletta — clothes peg; hair clip}

\label{lesson:IT-C531-molletta}



\begin{quote}
Before the new one: say the Italian for to hang, then the Italian for clothes horse.
\end{quote}

\subsection*{You'll want to know: la molletta}
\textbf{la molletta} — \textquotedblleft{}clothes peg; hair clip\textquotedblright{}.

\begin{quote}
Mi servono altre mollette per le lenzuola.
\end{quote}

\begin{quote}
\emph{I need more pegs for the sheets.}
\end{quote}

Also a hair clip.

\begin{cousinweb}
\emph{molletta} is a little \emph{molla} \textquotedblleft{}spring\textquotedblright{} (EtymDB), from \emph{mollare} \textquotedblleft{}to let go\textquotedblright{}, built on Latin \emph{mollis} \textquotedblleft{}soft\textquotedblright{} (Lewis \& Short).
\end{cousinweb}

\begin{grammarlens}[title={What you've built}]
One more word for telling, explaining and describing at B1.
\end{grammarlens}

\subsection*{Your turn}
\begin{itemize}
\raggedright
  \item \emph{Say it:} \emph{la molletta}
  \item \emph{Say it:} \emph{la molletta}, once more
  \item \emph{Say it:} a sentence of your own with \emph{la molletta}
  \item \emph{Recall:} say the Italian for to hang, then the Italian for clothes horse, then say \emph{la molletta} again
\end{itemize}

\begin{tcolorbox}[breakable,colback=teal!4,colframe=teal!35!black,title={Before you move on}]
What does la molletta mean? (\textquotedblleft{}Clothes peg; hair clip\textquotedblright{}.) Say it once more.
\end{tcolorbox}
\section[il ferro da stiro]{il ferro da stiro — iron (for clothes)}

\label{lesson:IT-C531-ferro-da-stiro}



\begin{quote}
Before the new one: say the Italian for clothes horse, then the Italian for clothes peg.
\end{quote}

\subsection*{You'll want to know: il ferro da stiro}
\textbf{il ferro da stiro} — \textquotedblleft{}iron (for clothes)\textquotedblright{}.

\begin{quote}
Attenta, il ferro da stiro è ancora caldo.
\end{quote}

\begin{quote}
\emph{Careful, the iron is still hot.}
\end{quote}

\emph{Stirare} — to iron.

\begin{cousinweb}
\emph{ferro da stiro} is built from \emph{ferro} \textquotedblleft{}iron\textquotedblright{}, inherited from Latin \emph{ferrum} (EtymDB), and \emph{stiro}, from \emph{stirare} \textquotedblleft{}to stretch, to iron\textquotedblright{}.
\end{cousinweb}

\begin{grammarlens}[title={What you've built}]
One more word for telling, explaining and describing at B1.
\end{grammarlens}

\subsection*{Your turn}
\begin{itemize}
\raggedright
  \item \emph{Say it:} \emph{il ferro da stiro}
  \item \emph{Say it:} \emph{il ferro da stiro}, once more
  \item \emph{Say it:} a sentence of your own with \emph{il ferro da stiro}
  \item \emph{Recall:} say the Italian for clothes horse, then the Italian for clothes peg, then say \emph{il ferro da stiro} again
\end{itemize}

\begin{tcolorbox}[breakable,colback=teal!4,colframe=teal!35!black,title={Before you move on}]
What does il ferro da stiro mean? (\textquotedblleft{}Iron (for clothes)\textquotedblright{}.) Say it once more.
\end{tcolorbox}
